% ===================================================================
%  કોષ્ટક નં. ૧૬ — મોટી દુકાન. --- બજાર. રમકડાં.
%
%  Traduction, faite en 2026, en gujarati d'aujourd'hui, sur l'ido de
%  texto/io. Regles de la colonne : voir 10-kostak-01.tex. Dernier
%  tableau du livret. Une seule numerotation. 34 blocs, 106 renvois.
%
%  TROISIEME ECART DECLARE, ET C'EST LE DERNIER DU RELEVE. Le bloc
%  t16-15-1 est aveugle : le fac-simile ido y compose
%  « \textsuperscript{d)} », sans parenthese ouvrante, et renvoji.py
%  ne voit donc que « 85 a b c e f 86 87 88 89 ». Un ecart declare ne
%  rend pas le bloc tolerant, il le rend INVISIBLE — le controle n'y
%  regarde plus rien. L'ordre juste est 85, a, b, c, d, e, f, 86, 87,
%  88, 89 ; il a ete pose a la main et recompte a la main, comme dans
%  les colonnes persane, haoussa, et comme au tableau 8 de celle-ci.
%
%  ET CE BLOC AVEUGLE EST JUSTEMENT CELUI QUI DIT LE PLUS DE CETTE
%  LANGUE. L'enfant regarde un tableau d'astronomie et declare n'y
%  rien comprendre. Or c'est le seul endroit du livret ou le gujarati
%  n'a pas un mot a chercher : ધૂમકેતુ la comete, ગ્રહણ l'eclipse,
%  ગોળાર્ધ l'hemisphere, ગ્રહ la planete, સૂર્યમંડળ le systeme solaire
%  — tous sanskrits, tous vieux de deux mille ans, tous encore dans le
%  journal du matin quand il y a une eclipse. LA SEULE CHOSE DE TOUTE
%  LA BOUTIQUE QUE L'ENFANT NE COMPREND PAS EST LA SEULE POUR
%  LAQUELLE CETTE LANGUE A LES MOTS LES PLUS ANCIENS.
%
%  Et les quatre points cardinaux (c) reviennent une troisieme fois.
%  Le persan, au meme renvoi, partageait sa rose des vents entre deux
%  langues, شمال et جنوب arabes contre خاور et باختر iraniens ; le
%  haoussa avait ses quatre mots a lui, arewa kudu gabas yamma, et
%  n'avait rien emprunte. Le gujarati fait le troisieme cas : ઉત્તર,
%  દક્ષિણ, પૂર્વ, પશ્ચિમ, quatre mots sanskrits d'un seul tenant, qui
%  servent a s'orienter dans une rue comme a orienter un temple. Trois
%  colonnes, un seul renvoi, trois reponses.
%
%  LE TITRE DU TABLEAU N'A PAS EU BESOIN D'ETRE TRADUIT. « La Bazaro »
%  est « બજાર », et c'est le seul des seize titres du livret dont les
%  deux langues possedaient deja le mot. Il est persan, et il est
%  arrive dans chacune par une route differente : dans le gujarati
%  directement, avec les marchands ; dans l'ido par le turc, l'italien
%  et le francais, apres un tour de la Mediterranee. Le livret dit
%  bazar a un lecteur gujarati qui dit bazar. C'est la derniere chose
%  que ce releve avait a trouver, et elle est dans le titre.
%
%  ONZIEME MOT PORTUGAIS, ET IL TOMBE AU DERNIER TABLEAU : « નાતાલ »
%  (81), Noel, de Natal. Le compte avait commence a la table et au
%  pretre, il passait par la pompe, l'hopital, la bonne d'enfants, la
%  pomme de terre, le chou et le pain ; il finit sous l'arbre de Noel
%  d'un grand magasin francais. ONZE MOTS EN SEIZE TABLEAUX, dont
%  aucun n'a l'air etranger a qui le parle.
%
%  Enfin la harpe (71) est « વીણા », et il faut une precaution. La
%  વીણા d'aujourd'hui est un luth a manche long ; mais celle qui est
%  sculptee a Sanchi et a Amaravati, il y a deux mille ans, est une
%  harpe arquee, et c'est ce mot-la qui a glisse d'un instrument a
%  l'autre. Le jouet pose sur l'etagere est ce que le mot voulait dire
%  avant. On l'ecrit sans hesiter : c'est la langue qui a bouge, pas
%  la chose.
% ===================================================================

%%K t16-apar-1 apar
\VUsaut{4.0mm}
\VUtitre{15.0pt}{60}{કોષ્ટક નં. ૧૬}
\VUsaut{2.0mm}
\VUpk{13.2pt}{40}{મોટી દુકાન. --- બજાર.}
\VUpk{13.2pt}{40}{રમકડાં.}
\VUsaut{3.4mm}

%%K t16-01-1 p
1. --- નવું વર્ષ પાસે આવે છે. મોટી દુકાનોએ \VUgras{રમકડાં}
\textsuperscript{(1)} અને નવા વર્ષની ભેટની પોતાની માંડણી તૈયાર કરી છે.

%%K t16-01-2 p
મેં મારા દાદાને કહ્યું કે એ સુંદર પ્રદર્શન જોવા મને બજારમાં લઈ જાય,
અને તેમણે હા પાડી.

%%K t16-02-1 p
2. --- પહેલાં અમે હથિયારની સુંદર ઢાલો આગળ થોભ્યા, જેમાં
\VUgras{બંદૂક}, \VUgras{લશ્કરી ટોપી}, \VUgras{વાંકી તલવાર},
\VUgras{તલવારનો પટ્ટો} અને \VUgras{ખભાની પટ્ટીઓની} જોડ છે, અથવા
\VUgras{ટોપ}, \VUgras{બખ્તર} \textsuperscript{(3)} અને
\VUgras{પિસ્તોલ} \textsuperscript{(4)}. એ બધામાં મને
\VUgras{સર્ચલાઇટ} \textsuperscript{(2)} કરતાં ઘણો વધારે રસ પડ્યો.

%%K t16-tit-1 sub
\VUsaut{3.0mm}
\VUpk{10.2pt}{40}{જોવા જેવાં દૃશ્યો.}
\VUsaut{2.6mm}

%%K t16-03-1 p
3. --- એ જ વિભાગમાં એક \VUgras{ચોકીદાર સિપાહી} \textsuperscript{(5)}
પોતાની \VUgras{ચોકીની ઓરડીમાં} હતો ; તે આખા પૂંઠાના પ્રાણીસંગ્રહ આગળ
ચોકી કરતો હોય એવું લાગતું હતું. ત્યાં કેટલાં બધાં પ્રાણી હતાં :
પોતાની દાંડી પર \VUgras{પોપટ} \textsuperscript{(6)}, નાક પર શિંગડાવાળો
\VUgras{ગેંડો} \textsuperscript{(7)}, \VUgras{રેન્ડિયર}
\textsuperscript{(8)}, \VUgras{શાહમૃગ} \textsuperscript{(9)}, અખરોટ
ફોડતો \VUgras{વાંદરો} \textsuperscript{(10)}, મરઘી ઉપાડી જતું
\VUgras{શિયાળ} \textsuperscript{(11)}, પ્રચંડ જડબાં ફાડતો
\VUgras{મગર} \textsuperscript{(12)}, \VUgras{ઊંટ} \textsuperscript{(13)},
નમણો \VUgras{ગ્રેહાઉન્ડ} \textsuperscript{(14)}, \VUgras{દીપડો}
\textsuperscript{(15)}, પછી બે પ્રાણી જેમને હું ઓળખતો ન હતો અને જે,
કહે છે તેમ, \VUgras{વીઝલ} \textsuperscript{(16)} અને \VUgras{ઝરખ}
\textsuperscript{(17)} છે. ત્યાં \VUgras{ચિત્તો} \textsuperscript{(18)}
અને \VUgras{વરુ} \textsuperscript{(21)} પણ હતાં ; એકદમ પાસે જ સુંદર
\VUgras{નાના ઉંદર} \textsuperscript{(19)} અને રબરની બહુ નાની
\VUgras{ઉંદરડીઓ} \textsuperscript{(20)} હતી. છેલ્લે
\VUgras{હિપોપોટેમસ} \textsuperscript{(22)} અને ખૂંધવાળા \VUgras{એક
ખૂંધના ઊંટે} \textsuperscript{(23)} સંગ્રહ પૂરો કર્યો. હું ત્યાં
કેટલાય કલાક રહ્યો હોત, જો પાસે જ \VUgras{સીસાના સિપાહીઓથી}
\textsuperscript{(29)} ભરેલી ભવ્ય છાવણી ન હોત, જેણે મારું ધ્યાન ખેંચ્યું.

%%K t16-tit-2 sub
\VUsaut{4.0mm}
\VUpk{10.2pt}{40}{સિપાહી અને લડાઈ.}
\VUsaut{2.8mm}

%%K t16-04-1 p
4. --- મારા દાદા, જે \VUgras{અપંગ} \textsuperscript{(24)} છે અને જેમને
એક \VUgras{ઘાને} લીધે લશ્કર છોડવું પડ્યું — એ જ ઘાએ તેમને
\VUgras{લશ્કરી ચાંદ} અપાવ્યો — તેમને મને એ \VUgras{લડાઈઓ} વિશે કહેવું
ખૂબ ગમે છે જેમાં તેમણે ભાગ લીધો હતો. તે નાના લઘુચિત્ર લશ્કરની જુદી
જુદી હિલચાલ સમજાવવામાં ખુશ હતા. બજાર જોવા આવેલા સાચા સિપાહીઓનું એક
જૂથ — એક \VUgras{ઇજનેરી સિપાહી} \textsuperscript{(25)}, માથે
\VUgras{બેરે ટોપીવાળો} \VUgras{આલ્પ્સનો શિકારી સિપાહી}
\textsuperscript{(26)}, પાયદળનો એક \VUgras{સાર્જન્ટ મેજર}
\textsuperscript{(27)} અને બાંય પર સોનેરી \VUgras{પટ્ટીવાળો}
\VUgras{તોપખાનાનો સાર્જન્ટ} \textsuperscript{(28)} — સમજૂતી સાંભળવા
અમારી પાસે થોભ્યું.

%%K t16-05-1 p
5. --- મારા દાદા, આટલા ચમકતા શ્રોતાઓ મળ્યા તેથી ખૂબ ગર્વ અનુભવીને,
તરત જ આખી લડાઈની યોજના ઊભી કરી દીધી.

%%K t16-05-2 p
કિલ્લેબંધ શહેર, પોતાના \VUgras{કોટ} અને \VUgras{ખાઈ} સાથે,
\VUgras{દુશ્મન લશ્કરે} ઘેર્યું હતું, જે લાંબા \VUgras{તોપમારા} પછી
\VUgras{હલ્લો કરવા} તૈયાર હતું. પણ \VUgras{ઘેરાયેલા}, ભૂખથી લાચાર
થઈને, \VUgras{ઘેરો ઘાલનારાની} \VUgras{છાવણી} પર \VUgras{બહાર ધસી જવા}
મથ્યા. \VUgras{તંબૂઓની} આસપાસ હઠીલી \VUgras{લડાઈથી} \VUgras{હુમલો}
પાછો હઠાવ્યા પછી, \VUgras{જીતેલા} ઘેરો ઘાલનારાએ \VUgras{હારેલા}
હુમલાખોરોનો પીછો કરવા પોતાની \VUgras{ઘોડેસવાર ટુકડીઓ} છોડી. તેમણે
\VUgras{પીછેહઠ કરી}, \VUgras{રણમેદાનને} \VUgras{મરેલા} અને
\VUgras{ઘાયલોથી} ઢાંકતાં ઢાંકતાં. \VUgras{સ્ટ્રેચરવાળા} એ ઘાયલોને
\VUgras{દવાખાનાની ગાડીમાં} લઈ ગયા, \VUgras{સર્જન} પાસે સારવાર કરાવવા.

%%K t16-05-3 p
\VUgras{ચોરસ} રચનારી કેટલીક \VUgras{બટાલિયનના} વીરોચિત સામના છતાં
\VUgras{હાર} પૂરી હતી, બાકીનું લશ્કર \VUgras{ભાગી ગયું}, પાછળ ઘણા
\VUgras{કેદીઓ} છોડીને. દુશ્મન શહેર કબજે કરીને પોતાની \VUgras{જીત}
પૂરી કરવા ગયો. કદાચ એ લડાઈનો અંત આવશે, અને \VUgras{યુદ્ધવિરામ} પછી
\VUgras{શાંતિ કરાર} પર સહી થઈ શકશે.

%%K t16-tit-3 sub
\VUsaut{4.4mm}
\VUpk{10.2pt}{40}{નવા વર્ષની ભેટ.}
\VUsaut{2.8mm}

%%K t16-06-1 p
6. --- જુવાન સિપાહીઓ, જે એ અનુભવીની રંગીન વાત સાંભળીને હસતા હતા,
તેમની પાસે બીજી થોડી સમજૂતી માગી. મારે તો, હું સુંદર \VUgras{ક્રૉસબો}
\textsuperscript{(32)} અને \VUgras{ધનુષ} \textsuperscript{(33)} તથા
તેના \VUgras{બાણને} \textsuperscript{(31)} જોવા દુકાનમાં ફર્યો. ત્યાં
મને પ્રાણીઓનો બીજો સંગ્રહ મળ્યો : બે \VUgras{ગાનારાં પક્ષી}
\textsuperscript{(34)} — યાંત્રિક \VUgras{બુલબુલ} અને ગળું ફાડીને
ગાતું \VUgras{વૉર્બ્લર} — અને વિચિત્ર પ્રાણીઓની હાર :
\VUgras{ચામાચીડિયું} \textsuperscript{(35)}, \VUgras{અજગર}
\textsuperscript{(36)}, \VUgras{કાચબો} \textsuperscript{(37)},
\VUgras{સીલ} \textsuperscript{(38)}, \VUgras{વ્હેલ}
\textsuperscript{(39)} અને ફુગાવેલો \VUgras{શાર્ક}
\textsuperscript{(40)}.

%%K t16-07-1 p
7. --- હું તેમને જોવા લાંબો સમય થોભ્યો નહીં, કેમ કે મેં જેનું સપનું
જોયું હતું તે જોયું : ભવ્ય \VUgras{સજાવટ} અને મોહક લાલ
\VUgras{પડદાવાળું} અતિ સુંદર \VUgras{કઠપૂતળીનું નાટકઘર}
\textsuperscript{(41)}. \VUgras{રંગમંચ} પર બે નટ કોઈ ખૂબ રસપ્રદ
\VUgras{નાટકમાં} \VUgras{ભૂમિકા} ભજવતા હોય એવું લાગતું હતું, કેમ કે
ખાનગી બેઠકોમાં ગોઠવેલા કે \VUgras{આરામખુરશીઓ} પર અને
\VUgras{વાદ્યવૃંદના નાયકની} પાછળ \VUgras{નીચેની બેઠકોમાં} બેઠેલા
પૂંઠાના નાના \VUgras{પ્રેક્ષકો} જોરથી તાળીઓ પાડતા હતા.

%%K t16-07-2 p
અફસોસ કે એ નાટકઘર બહુ મોંઘું હતું.

%%K t16-08-1 p
8. --- વળી, રમકડાં પસંદ કરવાં મુશ્કેલ હતું, કેમ કે શોકેસમાં બધી
જાતનાં રમકડાં ખડકાયેલાં હતાં : \VUgras{આર્લેકિન} \textsuperscript{(42)},
\VUgras{પુલ્ચિનેલો} \textsuperscript{(43)}, \VUgras{પિયેરો}
\textsuperscript{(44)}, પાંખો ફફડાવતાં \VUgras{ઘુવડ}
\textsuperscript{(45)}, સીટી વગાડતાં \VUgras{બ્લૅકબર્ડ}
\textsuperscript{(46)}, ભસતા \VUgras{પૂડલ કૂતરા},
\VUgras{ગ્રામોફોન} \textsuperscript{(47)} અને \VUgras{ધીરજની રમતો}.
એમાંના એક રમકડાએ મને ખૂબ મજા કરાવી : તે એક \VUgras{દરિયાઈ લડાઈ}
\textsuperscript{(48)} બતાવતું હતું, જેમાં એક \VUgras{વહાણ} પર
\VUgras{ચાંચિયા} હુમલો કરે છે, \VUgras{નાળિયેરીનાં ઝાડ} વાવેલા એક
દેશની પાસે.

%%K t16-noto-1 noto
\VUnotes{143.2mm}{%
\textit{(*) કોષ્ટક નં. ૬ જુઓ.}%
}

%%K t16-09-1 p
9. --- હું એ બધી અજાયબીઓ ખૂબ નિરાંતે જોઈ શક્યો, કેમ કે દુકાનમાં થોડા
જ માણસો હતા — માંડ થોડા લટાર મારનારા, એક \VUgras{ડ્રૅગૂન સિપાહી}
\textsuperscript{(49)} અને એક \VUgras{ઉઘરાણીદાર}
\textsuperscript{(50)}, જે મુખ્ય \VUgras{ગલ્લે} \textsuperscript{(51)}
\VUgras{હૂંડી} રજૂ કરતો હતો.

%%K t16-10-1 p
10. --- મેં મારા દાદાને પૂછ્યું કે \VUgras{કૅશિયરની} પાછળ પાટિયાં પર
મૂકેલી ચોપડીઓ શા કામની છે ; તેમણે મને જવાબ આપ્યો કે એ
\VUgras{ચોપડા} છે.

%%K t16-tit-4 sub
\VUsaut{3.6mm}
\VUpk{10.2pt}{40}{દુકાનની આસપાસ તાકવું.}
\VUsaut{2.8mm}

%%K t16-11-1 p
11. --- રમકડાંનો વિભાગ છોડીને અમે પલાણની દુકાને ગયા, \VUgras{પલાણ}
\textsuperscript{(53)}, \VUgras{પેંગડાં} \textsuperscript{(54)},
\VUgras{લગામ} \textsuperscript{(55)}, \VUgras{ચોકડું}
\textsuperscript{(56)} અને \VUgras{ચાબુક} પર નજર નાખવા. જ્યારે
\VUgras{વિભાગનો ઉપરી} \textsuperscript{(57)} મારા દાદાને થોડી માહિતી
આપતો હતો, ત્યારે હું \VUgras{ચિનાઈ માટીનાં વાસણ} અને
\VUgras{ક્રિસ્ટલની} \textsuperscript{(58)} માંડણી જોવા ગયો, જ્યાં
સુંદર \VUgras{સલાડનાં વાસણ} અને નાજુક \VUgras{ચાની કીટલીઓ}
\textsuperscript{(59)} છે.

%%K t16-12-1 p
12. --- પહેલા માળે જવા માટે અમે \VUgras{કૉર્સેટના}
\textsuperscript{(60)} શોકેસની પાછળથી પસાર થયા, મોજાંની દુકાન પાસે,
જ્યાં ખૂબ સુંદર \VUgras{અંદરનાં કપડાં} \textsuperscript{(63)} માંડેલાં
હતાં. પાસે જ એક \VUgras{વેચનારી બહેન} \textsuperscript{(61)} એક
\VUgras{ખરીદનારી બહેન} \textsuperscript{(62)} પાસે સુંદર
\VUgras{રેશમી કાપડ} \textsuperscript{(64)} પસંદ કરાવતી હતી.

%%K t16-12-2 p
દાદરે ચડતાં મેં જોયું કે દુકાન \VUgras{જાપાની ફાનસ}
\textsuperscript{(65)}, \VUgras{છત્રીઓ} \textsuperscript{(66)} અને
પ્રચંડ \VUgras{તાડનાં ઝાડથી} \textsuperscript{(70)} શણગારેલી છે.

%%K t16-13-1 p
13. --- પહેલા માળે બહુ જોવા જેવું ન હતું ; માત્ર રાચરચીલું (*),
\VUgras{તિજોરીઓ} \textsuperscript{(67)}, \VUgras{ખાનાંવાળી પેટીઓ}
\textsuperscript{(68)}, \VUgras{બાળગાડીઓ} \textsuperscript{(69)}. પણ
દુકાનનું આખું દૃશ્ય ખૂબ \VUgras{મજાનું} હતું.

%%K t16-14-1 p
14. --- અમારા પગ પાસે મેં વાજિંત્રો જોયાં : \VUgras{વીણા}
\textsuperscript{(71)}, \VUgras{ગિટાર} \textsuperscript{(72)},
\VUgras{મેન્ડોલિન} \textsuperscript{(73)}, \VUgras{ચાવીવાળાં કૉર્નેટ}
\textsuperscript{(74)}, \VUgras{ક્લેરિનેટ} \textsuperscript{(75)},
પોતાના \VUgras{ગજ} \textsuperscript{(76)} સાથે વાયોલિન અને
\VUgras{નગારું} \textsuperscript{(77)} સુધ્ધાં. થોડે આગળ, કાગળના
વિભાગ પાસે, એક \VUgras{વિદ્યાર્થી} \textsuperscript{(78)} એક
\VUgras{ગુમાસ્તા} \textsuperscript{(79)} પાસે નોટબુકની થેલીનો ભાવ
કરતો હતો. તેમની પાસે મેં એક સુંદર \VUgras{કક્કાની ચોપડી}
\textsuperscript{(80)} જોઈ.

%%K t16-14-2 p
દુકાનની વચ્ચે એક ઊંચું \VUgras{નાતાલનું ઝાડ} \textsuperscript{(81)}
ઊભું કરેલું હતું, જે \VUgras{મીણબત્તીઓ}, નાની નાની ચીજો અને બધી
જાતનાં રમકડાંથી ભરેલું હતું : \VUgras{લાકડાના ઘોડા}
\textsuperscript{(82)}, \VUgras{ઢીંગલીઓ} \textsuperscript{(83)}, વગેરે.

%%K t16-14-3 p
\VUgras{અત્તરની દુકાન} \textsuperscript{(84)}, પોતાની \VUgras{દાંતની
પાવડર} અને જુદાં જુદાં \VUgras{અત્તર} સાથે, ખૂબ સારી સુગંધ ફેલાવતી
હતી, જે અમારા સુધી પહોંચતી હતી.

%%K t16-tit-5 sub
\VUsaut{3.4mm}
\VUpk{10.2pt}{40}{અમે થોડું ખરીદીએ છીએ.}
\VUsaut{2.8mm}

%%K t16-15-1 p
15. --- મારા દાદાને સમય બહુ લાંબો લાગવા માંડ્યો, એટલે અમે મોટા
દાદરેથી નીચે ઊતર્યા. તે \VUgras{ખગોળના કોષ્ટક} \textsuperscript{(85)}
આગળ થોડું થોભ્યા, જે — તેમણે મને કહ્યું — \VUgras{ધૂમકેતુ}
\textsuperscript{(a)}, \VUgras{ચંદ્રગ્રહણ અને સૂર્યગ્રહણ}
\textsuperscript{(b)}, \VUgras{ચાર દિશાઓ} \textsuperscript{(c)}
\VUgras{(ઉત્તર, દક્ષિણ, પૂર્વ અને પશ્ચિમ)} વાળું દિશાચક્ર, બે
\VUgras{ગોળાર્ધનો} \textsuperscript{(d)} નકશો, \VUgras{ગ્રહો}
\textsuperscript{(e)} અને \VUgras{સૂર્યમંડળ} \textsuperscript{(f)}
બતાવે છે. એ બધામાં હું કશું સમજ્યો નહીં, અને મને તો પસાર થતા
\VUgras{માલ પહોંચાડનારા છોકરાની} \textsuperscript{(86)} પટ્ટીવાળી
વરદીમાં ઘણો વધારે રસ પડ્યો, અને મારી પાસે પડેલા સુંદર \VUgras{પંખામાં}
\textsuperscript{(87)}, જે \VUgras{પાકીટ} \textsuperscript{(88)} અને
\VUgras{ચામડાની થેલીઓ} \textsuperscript{(89)} પાસે હતા.

%%K t16-16-1 p
16. --- માંડણી પર મેં \VUgras{પીંછાં} \textsuperscript{(90)} અને
\VUgras{પટ્ટાનાં બકલ} \textsuperscript{(91)} પણ જોયાં, અને ટોપલાઓમાં
નમણી રીતે ગોઠવેલાં સુંદર \VUgras{કૃત્રિમ ફૂલ} \textsuperscript{(94)}.
\VUgras{વાયોલેટ}, \VUgras{ગિલિફ્લાવર}, \VUgras{હાયસિન્થ}
\textsuperscript{(95)}, \VUgras{જિરેનિયમ} \textsuperscript{(96)},
\VUgras{લિલી} \textsuperscript{(97)} અને \VUgras{નાસ્ટર્શિયમ}
\textsuperscript{(98)} ખૂબ સારી રીતે ઉતારેલાં હતાં.

%%K t16-tit-6 sub
\VUsaut{4.4mm}
\VUpk{10.2pt}{40}{મારા દાદાની ભેટ.}
\VUsaut{2.8mm}

%%K t16-17-1 p
17. --- મેં મારા દાદાને વિનંતી કરી કે \VUgras{દાંતના બ્રશમાંથી}
\textsuperscript{(92)} એક મારે માટે ખરીદે, જે \VUgras{વાળની પિન}
\textsuperscript{(93)} પાસેની પેટીમાં હતાં. તેમણે હા પાડી અને હું
પોતે જ ગલ્લે મારી ખરીદીના પૈસા ભરવા ગયો, જ્યાં પાસે જ એક
મોટું \VUgras{રવાનગીનું પોટલું} \textsuperscript{(100)} તૈયાર થતું
હતું, એક \VUgras{ગ્રાહક} \textsuperscript{(99)} માટે.

%%K t16-18-1 p
18. --- એકાએક કલાત્મક \VUgras{કાંસાની કલાકૃતિઓ}
\textsuperscript{(101)} પાસે થોભેલા લોકોમાં થોડી ધમાલ થઈ. એક
\VUgras{ચોર} \textsuperscript{(102)} હમણાં જ એક \VUgras{ઇન્સ્પેક્ટરે}
\textsuperscript{(103)} રંગેહાથ પકડ્યો હતો, જ્યારે તે કોઈને લૂંટવાની
કોશિશ કરતો હતો. તેને \VUgras{પકડવા} હવાલદારને બોલાવવાની કાળજી લેવાઈ,
અને અમે બહાર નીકળ્યા એ જ ઘડીએ ગુનેગારને કેદખાનામાં લઈ જવાતો હતો.

\VUsaut{6.0mm}
\VUfilet{22mm}
